% !TeX root = ../main.tex
\chapter{Related Work}
\label{chap:related-work}

\section{Modelling social norms and social influence}
\label{sec:related-norm-modelling}

Social norms enter environmental models through different assumptions about what individuals observe,
what they consider appropriate and how they respond to others. Prawitz et al.~\cite{prawitz2026review}
organise these representations in a systematic review of 153 studies on social norms in
social--ecological systems. They distinguish norms based on observed behaviour from norms based on
beliefs about what others approve of. They also distinguish how a norm affects decisions: through
imitation, anticipated approval or a direct influence on behavioural intention. This separates the origin
of normative information from its application in a decision rule. Within this sample, agent-based models
are the most common approach, followed by evolutionary game theory. Norms derived from observed behaviour
dominate, while interactions between norm types and responses to changing behavioural prevalence receive
less attention. The review also identifies limited empirical quantification of norm mechanisms, which
complicates the interpretation of model parameters. These findings make the choice of a norm
representation a substantive part of modelling human behaviour. The broader social-influence literature
illustrates why this choice matters. Flache et al.~\cite{flache2017influence} distinguish assimilative
influence, similarity-biased influence and repulsive influence. Individuals may move towards the opinions
of others, selectively approach similar individuals or move away from dissimilar individuals. Depending
on these assumptions and the interaction structure, models can produce consensus, separate opinion
clusters or polarisation. These outcomes concern distributions of individual opinions; they cannot be
inferred directly from a single aggregate mitigation share. Nevertheless, the comparison demonstrates how
alternative interaction rules can generate different collective dynamics. M\"uller-Hansen et
al.~\cite{muellerhansen2017behaviour} extend this modelling question to Earth system models by
distinguishing individual decision rules, interactions and approaches to aggregation. Agent-based
representations can retain differences between individuals and their contacts, whereas aggregate models
describe population-level quantities under assumptions about how individual decisions combine. The
appropriate representation therefore depends on which mechanisms the research question requires.
Together, these reviews motivate comparing the information that produces a norm signal and the decision
rule through which it changes behaviour.

\section{Concrete implementations and interacting norms}
\label{sec:related-norm-implementations}

Resource-use models provide concrete examples of how social influence affects environmental outcomes.
Sigdel et al.~\cite{sigdel2019convergence} couple the choice between sustainable and unsustainable
harvesting to social ostracism and three different ecological resource models. Under strong coupling,
otherwise different ecological systems can develop similar dynamics, and high social-learning rates can
generate stable cycles. The result shows that the speed of behavioural adjustment can change the
qualitative dynamics of the coupled system. Household models make the underlying interactions more
explicit. James and Rosenberg~\cite{james2022water} represent water use in an agent-based model informed
by household data from Logan, Utah. Their peer-pressure mechanism compares household consumption with the
median consumption in the household's social network. In the simulated interventions, peer pressure saved
more water than encouragement by a water manager, while combining the two strategies produced the largest
savings. Social influence thus depends on the information available to households and on how that
information enters their decisions. A separate issue is whether a change in attitude produces a
corresponding change in action. Torren-Peraire et al.~\cite{torrenperaire2024cultural} represent cultural
identity, social learning and individual barriers to several low-carbon behaviours. These barriers allow
attitudes and realised behaviour to differ. Their simulations show that cultural connections and
confirmation bias affect the spread of low-carbon behaviour, while a minority of green influencers alone
produces limited emissions reductions. The study also uses Sobol sensitivity analysis to examine how
uncertain model assumptions contribute to variation in the results. Bechthold et
al.~\cite{bechthold2025norms} distinguish descriptive norms, concerning observed behaviour, from
injunctive norms, concerning socially approved behaviour, in a model of renewable resource use. Both
norms evolve within a social network and interact with resource dynamics. Opposing normative pressures
can support intermediate outcomes, while group inertia and network structure influence the transition
between sustainable and unsustainable resource use. Here, norms change endogenously as part of the model;
this is distinct from a decision rule that explicitly responds to an observed behavioural trend. Bhalla
et al.~\cite{bhalla2024groundwater} further connect social and institutional norms with trust and
governance in a model of Tunisian groundwater-user groups. Norms could delay financial collapse in their
simulations but did not prevent hydrogeological collapse under the investigated conditions. They analyse
simulation outcomes using XGBoost, an ensemble of boosted decision trees. Across these studies, the
consequences of social influence depend on interactions with individual barriers, competing norms and
environmental constraints. A stronger norm can therefore have different effects depending on the
behaviour it supports and the feedbacks through which that behaviour affects the resource.

\section{Dynamic norm signals and perceived support}
\label{sec:related-dynamic-norms}

Information about behavioural change can affect decisions even when the current majority follows a
different practice. Sparkman and Walton~\cite{sparkman2017dynamic} investigate this mechanism in five
experiments involving sustainable behaviour. Messages indicating that behaviour was changing encouraged
movement towards the emerging practice, including when that practice remained a minority behaviour. Their
findings link this response to expectations about future prevalence and inferences about the importance
of the behaviour to others. Mortensen et al.~\cite{mortensen2019trending} similarly compare growing
minority norms with static minority information and control conditions. They find greater conformity to a
growing minority, with expectations of continued growth helping to explain the effect. Both studies
therefore support distinguishing how common a behaviour currently is from whether it is becoming more
common. They motivate the inclusion of trend information in a model, but do not establish a unique
mathematical trend function or a particular perception delay. Perception also matters when the relevant
signal concerns support for climate action. In a survey across 125 countries, Andre et
al.~\cite{andre2024support} find widespread support for climate action alongside systematic
underestimation of other people's willingness to contribute. This provides evidence that perceived
support can differ from reported support. The survey measures willingness and support, however, so its
percentages cannot be used directly as observed mitigation shares. A modelling application of dynamic
norm signals is provided by \'Alvarez-Antelo et al.~\cite{alvarez2025mobility}, who study shifts from air
to rail travel using system dynamics. Their dynamic norm mechanism responds to the increase in rail use
relative to its initial level, with a normalisation based on the initial air-travel share. It therefore
represents accumulated change from a reference situation rather than a recent rate of change. Their
scenarios show how personal environmental values and combinations of policy measures can affect the speed
of behavioural change and the resulting emissions. Read together, these studies motivate separate
treatment of current prevalence, perceived support and behavioural change. They also leave a modelling
question: how do different mathematical representations of these signals affect the dynamics when
behaviour and environmental conditions influence each other?

\section{Coupled behaviour--climate models}
\label{sec:related-coupled-models}

Coupled socio-climate models allow mitigation behaviour to change emissions and allow the resulting
climate conditions to influence later decisions. Beckage et al.~\cite{beckage2022incorporating} organise
relevant behavioural processes into cognition, contagion and response: how individuals perceive
information, how information and behaviour spread socially, and how individuals act. This framework helps
distinguish mechanisms that can otherwise be combined within a single behavioural response. In an earlier
model, Beckage et al.~\cite{beckage2018linking} link perceived risks from extreme events and a
representation of planned behaviour to the C-ROADS climate model. Behavioural feedbacks can increase or
decrease projected warming, depending on the assumed responses. Within their comparison, uncertainty
associated with behaviour is comparable in magnitude to uncertainty associated with the physical climate
system. The form of the behavioural response is therefore consequential for climate projections. Bury et
al.~\cite{bury2019pathways} provide the direct foundation for the aggregate socio-climate model used in
this thesis. Their model couples imitation between mitigating and non-mitigating individuals with carbon
and temperature dynamics. Individuals compare mitigation costs with perceived climate costs, while a
social-norm term reinforces the behaviour of the majority. Such reinforcement can initially impede
mitigation when mitigating individuals form a minority, even though the same mechanism supports
mitigation once it becomes sufficiently widespread. Social learning strongly affects the resulting
temperature trajectories. In the original model, perceived climate costs depend on a linear extrapolation
of recent warming, so expectations about future temperature also affect the timing of mitigation. This
assumption is relevant when interpreting adaptations of the model: responding to the current temperature
and responding to an extrapolated future temperature are different behavioural rules. Menard et
al.~\cite{menard2021conflicts} extend this line of work by introducing inequality and social groups with
different resources. Homophily, the tendency to interact more with similar individuals, changes how
mitigation strategies spread between groups. Greater inequality and homophily increased peak warming in
their simulations, and climate change could reinforce social polarisation. This extension demonstrates
that a population's internal structure can affect outcomes even when the model retains an
evolutionary-game approach to behavioural change.

More detailed coupled models examine particular consumption decisions or additional social and
institutional feedbacks. Eker et al.~\cite{eker2019diets} link endogenous dietary change to the FeliX
integrated assessment model. Their sensitivity analysis identifies social-norm effects and self-efficacy,
meaning the perceived ability to change behaviour, as major drivers of widespread dietary change. They
combine sensitivity analysis with scenario discovery to identify conditions under which dietary
transitions occur. Moore et al.~\cite{moore2022pathways} broaden the coupled system to include public
opinion, political decisions, technological change and climate feedbacks. Their ensemble of emissions
pathways depends strongly on climate perceptions, the cost and effectiveness of low-carbon technologies,
and the responsiveness of political institutions. The comparison with simpler behavioural models is
therefore partly a comparison of system boundaries: a mitigation response may be represented directly in
one model and mediated by policy or technology in another. Rajah et al.~\cite{rajah2025frida} investigate
endogenous dietary change in the fully coupled FRIDA model. Their framework includes personal norms,
descriptive norms, perceived accessibility and past behaviour. Perceptions and reference levels adjust
over time, allowing the behavioural response to depend on both the current situation and its
interpretation against previous conditions. Compared with demand projections driven by gross domestic
product, the behavioural framework produces lower projected food demand, while some scenarios include
later reversals. Its adjustment mechanisms provide a precedent for representing perception and reference
dynamics explicitly, although an adaptation must specify how these mechanisms are translated into its own
behavioural equations. A related approach to individual decisions and environmental feedbacks appears in
the InSEEDS preprint by Schwarz et al.~\cite{schwarz2025inseeds}. It couples farmer agents deciding
whether to adopt conservation tillage with the LPJmL vegetation model. Social networks, ecological
heterogeneity and behavioural inertia shape adoption, and the effects on soil carbon and yields depend on
the ecological context. The reported experiments hold the climate constant, so they examine coupled
human--biosphere processes without a changing climate feeding back into decisions. The study is
consequently useful for comparing heterogeneous adoption and environmental responses, while its scope
differs from models with endogenous global climate feedbacks. These examples show why similarities in
norm terminology do not imply that the surrounding behavioural and environmental mechanisms are
equivalent.

\section{Comparing formal representations}
\label{sec:related-model-comparison}

A comparison of norm models requires a way to describe their assumptions consistently. Schl\"uter et
al.~\cite{schlueter2017framework} address this problem through the Modelling Human Behaviour framework,
which maps behavioural theories to common elements of decision-making. It supports comparison of what
individuals perceive, how they evaluate alternatives and how they select an action. This makes
differences visible even when studies use different terminology or draw on different behavioural
theories. Muelder and Filatova~\cite{muelder2018formalizations} show that selecting a theory does not
uniquely determine its implementation. They compare alternative implementations of the theory of planned
behaviour in an agent-based model of household solar investment. Using the same theoretical foundation
and empirical data, different formal representations produced qualitatively and quantitatively different
outcomes. Their findings motivate explicit examination of decision rules and their mathematical
implementation alongside parameter uncertainty. The studies reviewed in this chapter provide evidence for
several norm mechanisms and demonstrate their importance in coupled environmental models. However, their
results come from different behavioural settings, environmental systems and assumptions about perception
and interaction. Comparing those results alone cannot establish how much a change in norm representation
alters the dynamics of a particular socio-climate model. This thesis addresses that question by comparing
selected norm representations within a shared climate-model framework. The comparison examines
behavioural trajectories, climate outcomes and parameter sensitivities, including the conditions
associated with qualitatively different simulation results. It includes both aggregate behavioural
equations and an agent-based decision rule, so the distinction between a change in the norm signal and a
change in the behavioural architecture remains explicit. Sobol analysis assesses the contribution of
parameter variation, while regression trees describe parameter combinations associated with different
outcomes. The purpose is to identify consequences of the specified representations within the
investigated model and parameter ranges. Such a comparison complements the existing mechanism-specific
studies without assuming that one formalisation provides a universally valid description of social norms.
